% ---
% id: meat-core/frank-wolfe/req:oracle-repair
% kind: req
% title: "Repair of an oracle return"
% summary: >-
%   An oracle return with a slightly negative spectrum may be rejected, or
%   projected onto the positive cone within a declared window of its own.
%   Every repair is recorded with its block, iterate and depth, apart from
%   iterate repairs, and the repaired return serves the search alone.
% prominence: standard
% status: implemented
% uses: [rem:fo-oracle-repair]
% mentions: [req:iterate-repair, req:verified-feasible-side, foundations/req:tolerance-scope]
% assumes: []
% refs: []
% provenance: newly-derived
% implements: certiqs_key.solvers.oracles.qics.adapter:QicsSectorChoiOracle
% lean:
% ---

No bracket reads a point the backend returned
(Requirement~\ref{req:verified-feasible-side}), so a repair of a return
changes only the search's next step and a recorded diagnostic. A window
bounds the displacement of the block but gives no bound on a pairing with
it (Remark~\ref{rem:fo-oracle-repair}), and the object repaired is a
\tref{computed object}, whose window is a number of its own
(\xref{foundations/req:tolerance-scope}\,(ii)).

\begin{requirement}\label{req:oracle-repair}
The engine may reject outright an oracle return a block of which has a
negative spectrum, or project that block onto the positive semidefinite cone
within the declared window \iface{oracle_repair_window}, a number of its own,
declared beside $\tnum$ and never as an instance of it. It records every
such repair, with the block, the iterate and the measured depth, in
\iface{optimizer_certificate}, kept apart there from the iterate repairs of
Requirement~\ref{req:iterate-repair}, and uses the repaired return only as
a step of the search and in the diagnostic pairing of
Requirement~\ref{req:verified-feasible-side}\,(vi), never in a bracket and
never as a bound on anything.
\end{requirement}
